\section{Method}
\label{sec:method}

\paragraph{Hardware and isolation.}
One NVIDIA GeForce RTX~4090 (24\,GB, 450\,W stock limit, PCIe~3.0~x16), an Intel Core
i9-10900K and 62\,GB of RAM, under Ubuntu with driver 580.178. Every engine is built and run in a
throw-away Docker container; the model files are mounted read-only. One server runs at a time.

\paragraph{Decode speed.}
A client streams each answer and computes decode speed as $(n-1)/(t_\text{last}-t_\text{first})$,
where $n$ is the number of generated tokens and $t$ the arrival times of the first and last
streamed tokens, so prompt processing is excluded and every engine is measured the same way. The
short benchmark has six prompts (two coding tasks, a short story, a technical explanation, a word
problem and a table), up to 1024 tokens each, with the sampling the model card recommends
(temperature 0.6, top-$p$ 0.95, top-$k$ 20, min-$p$ 0, presence penalty 0) and a fixed seed. We
report the mean over the prompts. Final configurations were repeated three times; with a fixed
seed the repetitions agree within 1\,tok/s, while separate server starts vary by about 5\%.
The long-prompt benchmark prefixes three questions with 29{,}005 tokens of CUDA source code.

\paragraph{Accuracy.}
GSM8K (first 250 test problems) and HumanEval
(all 164, pass@1, run against the official tests), greedy, thinking off.

\paragraph{Context capacity.}
For NInfer, the memory planner decides at startup whether a configuration fits, so we
binary-search the largest \texttt{--max-context} (in steps of 4096 tokens) that starts. For
OpenJet, which allocates at startup and fails with \texttt{cudaMalloc} out of memory, we do the
same. A start is not proof: each retained configuration is then filled to within 2--3\% of its
limit with real source code and queried (next paragraph); we call a context \emph{usable} only
if the server answers every question and is still alive afterwards.

\paragraph{KV precision.}
Two measurements. (i) Code perplexity: NInfer's offline evaluator scores 359{,}927 tokens of
CUDA/C++ source with windows of 64K (stride 32K) or 128K (stride 64K) tokens, for each KV format;
we patched its command line to accept the rotated and lattice formats that the server already
supports. For llama.cpp, \texttt{llama-perplexity} scores the same text in 64K chunks. For
azllm, \texttt{tqkld} first dumps the top-16 log-probabilities of the \texttt{q8} cache as a
reference, then scores the other formats against it (KL divergence, top-1 agreement and NLL),
on four 64K chunks of which the last 32K positions are scored. Within each engine, the most
precise format it offers is the reference. (ii) Long-context recall: the prompt holds up to 252K tokens of
llama.cpp source, then twelve questions ask the model to copy, exactly, the first 16 lines of
the body of a named function in a named file, at depths spread from 10\% to 98\%. A function
qualifies only if its name is defined once in its file. We report the share of lines copied
exactly.

\paragraph{Weight precision.}
The three weight files are compared on the same 98{,}304 positions: chunks 1--3 of 64K tokens,
last 32K of each, with 8-bit keys and values. \texttt{llama-perplexity} with 64K chunks,
\texttt{tqkld --tail 32768}, and NInfer's windows starting at 64K, 128K and 192K (stride 32K)
all score exactly these positions, and the three tokenizers produce the same 359{,}928 tokens. The
comparison includes each engine's own arithmetic, not the weights alone, and there is no BF16
reference: it tells which pipeline predicts the code best, not which is closest to the original
model.
